Put \(p_{a,n}(y,z)=q_{a,n}(y,z)r_n(z)\).  Products preserve the declared anisotropic H\"older bounds up to a constant depending only on \(C,L\).  At the pilot resolutions, the terminal-scale wavelet tails obey
\[
\max_a\|p_{a,n}-\Pi_{\bar J_y,\bar J_z}p_{a,n}\|_\infty
+\|r_n-\Pi_{\bar J_z}r_n\|_\infty
\lesssim2^{-s_y\bar J_y}+2^{-s_z\bar J_z}lesssim\varepsilon_{\mathrm{pil},n}. \tag{1}
\]
For a retained coefficient, exponential Markov's inequality applied to the centered bounded wavelet summands gives a tail bounded by \(2\exp\{-c m x^2/\sigma^2\}\) throughout the variance range used below.  At each point only a fixed number of compactly supported wavelets at a terminal scale are nonzero.  Taking a union bound over \(O(2^{\bar J_y+d\bar J_z})\) coefficients, and substituting the declared resolutions, therefore yields
\[
\max_{a,e}\|\widetilde p_{a,n}^{(-e)}-p_{a,n}\|_\infty
+\max_e\|\widetilde r_n^{(-e)}-r_n\|_\infty
=O_P(\varepsilon_{\mathrm{pil},n})                                        \tag{2}
\]
uniformly over the model class.  On the event where \(\widetilde r\ge c/2\), the identity
\[
\frac{\widetilde p}{\widetilde r}-\frac p r
=\frac{\widetilde p-p}{\widetilde r}
+p\frac{r-\widetilde r}{r\widetilde r}                                   \tag{3}
\]
and \(r\ge c\) give the same rate for each raw conditional pilot.  Since each true density lies inside the clipping interval \([c/2,2C]\), (2)--(3) imply \(P(\cap_e\mathcal E_n^{(-e)})\to1\).  Clipping and normalization are locally Lipschitz on this event, while by construction they enforce positivity and exact unit conditional mass everywhere.  This proves the first display and the clipping conclusions.

Let
\[
A_z=P_{J_y}\{-\partial_y(\bar q_n\partial_y)\}P_{J_y},
\qquad
\widehat A_z=P_{J_y}\{-\partial_y(\widehat{\bar q}_n^{(-e)}\partial_y)\}P_{J_y}.
\]
On the zero-mean subspace, the one-dimensional inequality used in the spectrum proof and ellipticity give \(A_z\succeq c_P I\) uniformly.  The weak forms show
\[
\|\widehat A_z-A_z\|_{H^1_0/\mathbb R\to H^{-1}}
\lesssim\|\widehat{\bar q}_n^{(-e)}-\bar q_n\|_\infty.
\]
The resolvent identity
\[
(\widehat A_z+\lambda_nI)^{-1}-A_z^{-1}
=(\widehat A_z+\lambda_nI)^{-1}
\{A_z-\widehat A_z-\lambda_nI\}A_z^{-1}                                 \tag{4}
\]
and the uniform coercivity bounds therefore give operator error \(O_P(\varepsilon_{\mathrm{pil},n}+\lambda_n)\).  Integrating the matrix entries against \(\widehat r_n\), and using (2), proves
\[
\max_e\|\widehat H_{J_y,J_z,n}^{(-e)}-H_{J_y,J_z,n}\|_{\mathrm{op}}
=O_P(\varepsilon_{\mathrm{pil},n}+\lambda_n).
\]

We next bound the deterministic population correction.  Parameterize the pair by \((\bar q,\Delta)\).  Exchanging the arms replaces \(\Delta\) by \(-\Delta\) and leaves squared transport unchanged, so the functional is even in \(\Delta\).  Twice differentiating the quantile identity used in the cosine proof, on the uniformly positive density set \([c/2,2C]\), shows that the equality Hessian is Lipschitz in \(\bar q\), and that its difference from the Hessian at \((\bar q,\Delta)\) is bounded by \(C\|\Delta\|_\infty\) as a bilinear form.

The required interpolation bound follows directly from the anisotropic H\"older constraint.  If \(|\Delta|\) attains amplitude \(A\), its H\"older bound forces amplitude at least a fixed multiple of \(A\) on an outcome interval of length \(A^{1/s_y}\) and a covariate rectangle of volume \(A^{d/s_z}\).  The zero outcome mass forces its antiderivative to attain order \(A^{1+1/s_y}\) on an outcome interval of length of order \(A^{1/s_y}\).  Ellipticity and the energy identity in the cosine proof thus give
\[
\rho_n^2\gtrsim A^{2+3/s_y+d/s_z},
\qquad
\|\Delta_n\|_\infty
\lesssim\rho_n^{\,2/(2+3/s_y+d/s_z)}.                                     \tag{5}
\]
Taylor's identity with integral remainder, (5), and the pilot rate give
\[
O_P\!\left\{\varepsilon_{\mathrm{pil},n}^3
+\rho_n^{\,2/(2+3/s_y+d/s_z)}\varepsilon_{\mathrm{pil},n}^2\right\}.      \tag{6}
\]
The first term is the equality third-order remainder and the second is the unequal-law Hessian mismatch.  Replacing the equality Hessian by the ridge pilot adds
\(O_P\{(\varepsilon_{\mathrm{pil},n}+\lambda_n)\varepsilon_{\mathrm{pil},n}^2\}\), which has the same order because \(\lambda_n=\varepsilon_{\mathrm{pil},n}\).  Sample splitting makes the remaining centered operator terms
\[
o_P(1)\left\{\frac{\rho_n}{\sqrt n}+\frac{A_{J_y,J_z,n}}n\right\}.
\]
This proves the stated pilot, ridge, and operator contribution and its negligibility criterion.

Finally, the spectrum proposition gives fixed lower and upper bounds for the feature covariance, while compact wavelet support gives \(\|\Phi(W)\|_2^2\lesssim D_{J_y,J_z}\).  Expanding the matrix exponential moment of each centered rank-one summand and taking a dimension union bound yields
\[
\max_e\|\widehat\Gamma_{J_y,J_z,n}^{(-e)}-\Gamma_{J_y,J_z,n}\|_{\mathrm{op}}
=O_P\!\left(\sqrt{\frac{D_{J_y,J_z}\log n}{m}}
+\frac{D_{J_y,J_z}\log n}{m}\right)=o_P(1).                              \tag{7}
\]
On the resulting fixed spectral interval, the square-root map is Lipschitz, as follows by applying the scalar integral representation of the positive square root to the two matrices.  Combining (4) and (7) in the definition of \(K\) proves
\[
\max_e\|\widehat K_{J_y,J_z,n}^{(-e)}-K_{J_y,J_z,n}\|_{\mathrm{op}}=o_P(1).
\]
The condition \(D_{J_y,J_z}\log n/m=o(1)\) was used only in (7), so it is needed only for the observed spectral proxy.